\documentclass[12pt]{amsart}

%% ─── Packages ────────────────────────────────────────────────────────────────
\usepackage{amsmath,amssymb,amsthm}
\usepackage{geometry}
\geometry{margin=1in}
\usepackage{hyperref}
\hypersetup{colorlinks=true,linkcolor=blue,citecolor=blue,urlcolor=blue}
\usepackage{booktabs}
\usepackage{enumitem}
\usepackage{microtype}

%% ─── Theorem Environments ────────────────────────────────────────────────────
\theoremstyle{plain}
\newtheorem{theorem}{Theorem}[section]
\newtheorem{proposition}[theorem]{Proposition}
\newtheorem{lemma}[theorem]{Lemma}
\newtheorem{corollary}[theorem]{Corollary}
\newtheorem{conjecture}[theorem]{Conjecture}

\theoremstyle{definition}
\newtheorem{definition}[theorem]{Definition}
\newtheorem{example}[theorem]{Example}

\theoremstyle{remark}
\newtheorem{remark}[theorem]{Remark}

%% ─── Notation ────────────────────────────────────────────────────────────────
\newcommand{\Z}{\mathbb{Z}}
\newcommand{\Q}{\mathbb{Q}}
\newcommand{\R}{\mathbb{R}}
\newcommand{\Ck}{\mathcal{C}_k}
\newcommand{\thetak}{\theta_k}
\newcommand{\rhoc}{\rho^{\mathrm{config}}}

%% ─── Metadata ────────────────────────────────────────────────────────────────
%% arXiv: primary math.NT, cross-list math.HO, cs.IT
\title[Constitutional Forcing]{Constitutional Forcing}

\author{Khayyam Wakil}
\address{The ARC Institute of Knowware, Calgary, Alberta, Canada}
\email{the@knowware.institute}
\date{February 2026}

\keywords{constitutional forcing, constitutional constraint, configuration
  density, forced governing constant, sieve theory, level of distribution,
  Shannon entropy, channel capacity, Kolmogorov scaling, $k$-hierarchy,
  Shannon--Wakil effect, cross-domain structure}

\subjclass[2020]{
  11N35,  % Sieves (primary)
  11N13,  % Primes in progressions
  94A15,  % Shannon theory
  76F02,  % Turbulence, fluid dynamics
  01A99   % Miscellaneous topics: history
}

%% ─── Abstract ────────────────────────────────────────────────────────────────
\begin{document}

\begin{abstract}
We introduce and formally define \emph{Constitutional Forcing}: a mechanism
by which the irreducible algebraic structure of a mathematical system uniquely
determines a governing constant, leaving no freedom in its value. The constant
is not fitted, not approximated, and not chosen---it is forced. Changing it
requires changing the structure of the system itself.

A system of depth $k$ has exactly $k$ independent \emph{constitutional
constraints}, operating on a binary configuration space of size $2^k$.
Each constraint eliminates exactly one configuration class, leaving $2^k - k$
valid configurations. The forced governing constant is the configuration density:
\begin{equation*}
  \thetak = \frac{2^k - k}{2^k}.
\end{equation*}

We exhibit three proven instances across three independent domains.
In \emph{prime arithmetic} ($k=1$, $k=3$): the Bombieri--Vinogradov level
of distribution $\theta_1 = 1/2$ and the cascade moduli level of distribution
$\theta_3 = 5/8$~\cite{WakilEisenstein,WakilW3} are both constitutional.
In \emph{information theory} ($k=1$): Shannon's entropy $H$ and channel
capacity $C$ are the forcing constants of the AEP and the Channel Coding
Theorem~\cite{Shannon1948}.
In \emph{fluid dynamics} ($k=2$): Kolmogorov's dissipation exponent
$\alpha = 3/4$ is forced by two incompressibility constraints on an
eight-dimensional triadic interaction space~\cite{Kolmogorov1941}.

We prove the Counting Theorem: exactly $k$ configurations are invalid in any
$k$-fold constitutionally constrained system. We verify that all three domains
satisfy the formal definition. We state the Constitutional Forcing Conjecture:
every system with $k$ independent constitutional constraints exhibits
$\thetak = (2^k - k)/2^k$ as its forcing constant.

This paper is the foundation of a seven-paper program. The remaining six papers
prove the instances, derive consequences, and explore the cross-domain
structure of the mechanism.
\end{abstract}

\maketitle
\tableofcontents

%% ─────────────────────────────────────────────────────────────────────────────
\section{What Constitutional Forcing Is}
%% ─────────────────────────────────────────────────────────────────────────────

\subsection{The problem this paper solves}

Mathematics is full of constants that appear to be chosen but turn out to be
forced. Euler's $e$, the golden ratio $\varphi$, Kolmogorov's $3/4$, Shannon's
channel capacity $C$, the Bombieri--Vinogradov exponent $1/2$, Pascadi's
$5/8$~\cite{Pascadi2025}: in each case, someone could in principle have asked
``why this value and not another?'' The answer, in every case, is that no other
value is consistent with the algebraic structure of the problem. The constant is
not discovered by fitting---it is forced by the geometry.

This paper names, defines, and formalises that forcing mechanism.

We call it \emph{Constitutional Forcing}. The name is deliberate. A
constitutional constraint is one that is not imposed from outside the system but
is \emph{inherent to its constitution}---to what the system fundamentally is.
Constitutional Forcing is what happens when a system's own constitution leaves
no room for any governing constant other than the one it has.

\subsection{Why this needs a paper}

The mechanism has been operating silently across mathematics for at least a
millennium---Khayyam's Triangle~\cite{WakilKhayyam} encodes it in the mod-3
fractal structure that waited nine centuries for the right question. Shannon
found two instances in 1948 without naming the mechanism. Kolmogorov found
another in 1941. The prime arithmetic instance was established in
2025--2026~\cite{Pascadi2025,WakilEisenstein}.

None of these works named what they had found. Each field had its own instance,
its own language, its own proof methods. The mechanism---the common cause---was
invisible because no one had looked across all three at once.

This paper looks across all three. It names the mechanism. It defines it
precisely enough that a referee can check whether a new system exhibits it.
And it positions the six companion papers as instances and consequences of the
same foundational definition.

\subsection{Relationship to the Shannon--Wakil Effect}

Constitutional Forcing is the \emph{mechanism}. The Shannon--Wakil
Effect~\cite{WakilShannon} is the \emph{observable}: the forced dimensional
reduction of an exponential configuration space that Constitutional Forcing
produces. The relationship is:
\[
  \text{Constitutional Forcing}
  \;\xrightarrow{\text{acts on}}\;
  \text{exponential space}
  \;\xrightarrow{\text{produces}}\;
  \text{Shannon--Wakil Effect}.
\]
Defining the Shannon--Wakil Effect without first defining Constitutional
Forcing is like defining combustion without defining oxygen. The present paper
provides the missing definition.

%% ─────────────────────────────────────────────────────────────────────────────
\section{The Definitions}
\label{sec:definitions}
%% ─────────────────────────────────────────────────────────────────────────────

\subsection{Constitutional constraints}

\begin{definition}[Constitutional constraint]
\label{def:constitutional}
Let $S$ be a mathematical system with a binary configuration space at each
operational level. A constraint $\mathcal{C}_i$ at level $i$ is
\emph{constitutional} if it satisfies all three of the following:
\begin{enumerate}[label=\textup{(C\arabic*)}]
  \item \emph{Intrinsic.} $\mathcal{C}_i$ is forced by the irreducible algebraic
    structure of $S$---it is not a free parameter, an external imposition, or
    a modelling choice.
  \item \emph{Independent.} $\mathcal{C}_i$ is structurally independent of
    constraints at other levels: it cannot be derived from, or reduced to, a
    combination of constraints at levels $j \neq i$.
  \item \emph{Binary-eliminating.} $\mathcal{C}_i$ eliminates exactly one
    configuration class at level $i$---the one where the constraint fails.
    It does not eliminate more (over-constrained) or fewer (under-constrained).
\end{enumerate}
\end{definition}

\begin{remark}
Condition (C1) is the most important and the hardest to fake. The test is
counterfactual: \emph{could one choose a different value without changing the
algebraic structure of the system?} If yes, the constraint is not constitutional.
If no---if changing the constraint requires changing what the system fundamentally
is---then it is.
\end{remark}

\begin{remark}
Condition (C2) is what distinguishes Constitutional Forcing from ordinary
constraint satisfaction. A system with $k$ coupled constraints (each derivable
from a single underlying constraint) has effective depth $1$, not depth $k$.
The Sophie Germain prime pairs are the canonical example: two modular conditions
that appear independent but are coupled through a single algebraic relationship,
giving effective depth $1$ and forcing constant $\theta = 1/2$, not the na\"ive
$\theta_2 = 1/2$ that the formula also predicts for $k=2$ (for independent
reasons)~\cite{WakilShannon}.
\end{remark}

\subsection{Constitutional Forcing}

\begin{definition}[Constitutional Forcing]
\label{def:CF}
A system $S$ exhibits \emph{Constitutional Forcing of depth $k$} if:
\begin{enumerate}[label=\textup{(CF\arabic*)}]
  \item \emph{Configuration space.} $S$ has a binary configuration space
    $\Ck = \{0,1\}^k$ of size $2^k$ at depth $k$. Each configuration
    $c = (c_1, \ldots, c_k)$ records whether the constraint at level $i$ is
    satisfied ($c_i = 1$) or violated ($c_i = 0$).
  \item \emph{$k$ independent constitutional constraints.} $S$ has exactly
    $k$ constitutional constraints $\mathcal{C}_1, \ldots, \mathcal{C}_k$
    in the sense of Definition~\ref{def:constitutional}, one per level.
  \item \emph{Forced configuration density.} The valid configurations---those
    not eliminated by any $\mathcal{C}_i$---number exactly $2^k - k$,
    giving forced governing constant
    \begin{equation}
    \label{eq:theta}
      \thetak = \frac{2^k - k}{2^k}.
    \end{equation}
  \item \emph{Uniqueness.} No other value of $\thetak$ is consistent with both
    the configuration space structure and the $k$ constitutional constraints.
    Changing $\thetak$ requires either removing a constitutional constraint
    (reducing depth) or adding one (increasing depth).
\end{enumerate}
\end{definition}

\begin{remark}
Condition (CF4) is what makes Constitutional Forcing a mechanism and not merely
a formula. The formula $(2^k - k)/2^k$ is a consequence, not a definition. The
definition is the forcing structure: the right-hand side is forced to equal the
left-hand side by (CF1)--(CF3), and (CF4) asserts that nothing else can equal it.
\end{remark}

\subsection{The Counting Theorem}

The following theorem establishes that (CF3) follows from (CF1) and (CF2): the
forced count $2^k - k$ is not an additional assumption but a consequence of
having exactly $k$ independent constitutional constraints.

\begin{theorem}[Counting Theorem]
\label{thm:count}
Let $S$ exhibit Constitutional Forcing of depth $k$. Then exactly $k$
configurations in $\Ck = \{0,1\}^k$ are invalid, and the forced governing
constant is $\thetak = (2^k - k)/2^k$.
\end{theorem}

\begin{proof}
By induction on $k$.

\textit{Base case} ($k = 1$). One level, configuration space $\{0,1\}$.
The single constitutional constraint eliminates the configuration where it fails
(exactly one, by (C3)). Invalid count: $1 = k$. Valid configurations:
$2 - 1 = 1$. Forced constant: $\theta_1 = 1/2$. \checkmark

\textit{Inductive step.} Assume the result for $k-1$: a system of depth $k-1$
has exactly $k-1$ invalid configurations out of $2^{k-1}$.

Passing to depth $k$, each configuration in $\mathcal{C}_{k-1}$ splits into
two in $\Ck$ (the new level either holds or fails). The previously invalid
$k-1$ configurations in $\mathcal{C}_{k-1}$ remain invalid (failing any
constraint is sufficient for invalidity). The new constitutional constraint
$\mathcal{C}_k$ is independent of the previous $k-1$ (by (C2)), so it
eliminates exactly one configuration that was valid at levels $1, \ldots, k-1$
but fails at level $k$ (by (C3)).

Total invalid: $(k-1) + 1 = k$. Valid: $2^k - k$. Forced constant:
$\thetak = (2^k - k)/2^k$.
\end{proof}

Table~\ref{tab:theta} displays the forced constants for $k = 1, \ldots, 6$.

\begin{table}[h]
\centering
\caption{Forced governing constants $\thetak = (2^k - k)/2^k$.}
\label{tab:theta}
\begin{tabular}{rrrrl}
\toprule
$k$ & $2^k$ & Invalid & Valid & $\thetak$ \\
\midrule
$1$ & $2$  & $1$ & $1$  & $1/2$ \\
$2$ & $4$  & $2$ & $2$  & $1/2$ \\
$3$ & $8$  & $3$ & $5$  & $5/8$ \\
$4$ & $16$ & $4$ & $12$ & $3/4$ \\
$5$ & $32$ & $5$ & $27$ & $27/32$ \\
$6$ & $64$ & $6$ & $58$ & $29/32$ \\
\bottomrule
\end{tabular}
\end{table}

%% ─────────────────────────────────────────────────────────────────────────────
\section{Three Instances}
\label{sec:instances}
%% ─────────────────────────────────────────────────────────────────────────────

We exhibit Constitutional Forcing in three independent mathematical domains.
Each instance is verified against Definition~\ref{def:CF}. The proofs of the
underlying results live in the companion papers; here we verify only that each
system satisfies the definition.

\subsection{Prime arithmetic: Bombieri--Vinogradov ($k = 1$)}

The Bombieri--Vinogradov theorem~\cite{BV1965} establishes that primes are
equidistributed in arithmetic progressions to modulus $q \leq x^{1/2}$. The
exponent $1/2$ is the level of distribution.

\begin{proposition}[BV is constitutionally forced]
\label{prop:BV}
The Bombieri--Vinogradov level of distribution $\theta_1 = 1/2$ is
constitutionally forced at depth $k = 1$.
\end{proposition}

\begin{proof}[Verification against Definition~\ref{def:CF}]
(CF1): Configuration space $\{0,1\}^1 = \{0,1\}$, size $2$.

(CF2): One constitutional constraint---primality itself. A positive integer
either is prime (constraint satisfied, $c_1 = 1$) or is not (constraint failed,
$c_1 = 0$). This is intrinsic (it is what primality means), independent (there
is only one constraint), and binary-eliminating.

(CF3): Invalid configurations: $1$ (the non-prime configuration). Valid: $1$.
Forced constant: $\theta_1 = 1/2$.

(CF4): The exponent $1/2$ cannot be improved unconditionally for all moduli
without additional structural input (such as the Generalised Riemann
Hypothesis). Any improvement requires either a new structural constraint
(increasing depth $k$) or a restriction to a structured family of moduli.
Both are exactly what Pascadi~\cite{Pascadi2025} and the companion
papers~\cite{WakilEisenstein} do: they increase depth to $k = 3$ by
restricting to cascade moduli.
\end{proof}

\subsection{Prime arithmetic: Cascade moduli ($k = 3$)}

For cascade moduli $q = 3^K$, the companion paper~\cite{WakilEisenstein}
establishes level of distribution $\theta_3 = 5/8$.

\begin{proposition}[Cascade moduli are constitutionally forced]
\label{prop:cascade}
The level of distribution $\theta_3 = 5/8$ for cascade moduli is
constitutionally forced at depth $k = 3$.
\end{proposition}

\begin{proof}[Verification against Definition~\ref{def:CF}]
(CF1): Configuration space $\{0,1\}^3$, size $8$.

(CF2): Three constitutional constraints. For a twin prime pair $(p, p+2)$
with $p > 3$:
\begin{itemize}
  \item Level $1$ ($\bmod\,3$): $p \equiv 2 \pmod{3}$ is forced---the only
    residue class allowing both $p$ and $p+2$ to avoid divisibility by $3$.
  \item Level $2$ ($\bmod\,9$): The survival condition at modulus $9$ is
    forced by the ramification $(3) = (1-\omega)^2$ in
    $\Z[\omega]$~\cite{WakilEisenstein}.
  \item Level $3$ ($\bmod\,27$): The survival condition at modulus $27$ is
    similarly forced by the three-level filtration of $\Z[\omega]/(3^K)$.
\end{itemize}
Each constraint is intrinsic to the twin prime condition and the Eisenstein
integer structure, independent of the others (verified via CRT
in~\cite{WakilRamanujan}), and binary-eliminating (each eliminates exactly one
of the eight configuration classes).

(CF3): Invalid configurations: $3$ (one per level). Valid: $8 - 3 = 5$.
Forced constant: $\theta_3 = 5/8$.

(CF4): The uniqueness of $5/8$ is proved in Theorem~5.1 of~\cite{WakilEisenstein}:
neither the Eisenstein filtration alone (giving ceiling $2/3$) nor the cyclic
group cancellation alone achieves $5/8$; both are required, and together they
force exactly $5/8$.
\end{proof}

\subsection{Information theory: Shannon's theorems ($k = 1$)}

\begin{proposition}[Shannon's constants are constitutionally forced]
\label{prop:shannon}
Shannon's entropy $H$ and channel capacity $C$~\cite{Shannon1948} are
constitutionally forced at depth $k = 1$ in their respective systems.
\end{proposition}

\begin{proof}[Verification against Definition~\ref{def:CF}]
\textit{AEP case.} The system is a memoryless source with alphabet $\mathcal{A}$
and distribution $p$.

(CF1): Configuration space is $\{0,1\}^1$ (sequences are either typical or not).

(CF2): One constitutional constraint---typicality. A sequence $x^n$ is typical
if $|-\frac{1}{n}\log p(x^n) - H| < \varepsilon$. This is intrinsic to the
source distribution $p$: it is not a free parameter.

(CF3): Invalid configurations: $1$ (the atypical class). Valid: $1$ (the typical
class). Forced constant: $\theta_1 = H / \log|\mathcal{A}|$---the fraction of
the full space occupied by the typical set.

(CF4): The entropy $H$ uniquely determines the size of the typical set. No other
value is consistent with $p$. Changing $H$ requires changing $p$, which changes
the system.

\textit{Channel Coding case.} Identical structure with $C = \max_{p(x)} I(X;Y)$
playing the role of $H$. The constitutional constraint is reliable
decodability---a rate either achieves it (below $C$) or does not (above $C$).
The capacity $C$ is forced by the channel $p(y|x)$; it cannot be changed
without changing the channel.
\end{proof}

\subsection{Fluid dynamics: Kolmogorov's $\alpha = 3/4$ ($k = 2$)}

Kolmogorov's 1941 theory~\cite{Kolmogorov1941} predicts that the vorticity
dissipation exponent satisfies $|\omega|_\infty \leq C \cdot
\varepsilon^{3/4}/\nu^{3/4}$, corresponding to scaling exponent $\alpha = 3/4$.
This exponent has been empirically validated for 84 years; what has been lacking
is an \textit{a priori} derivation.

\begin{proposition}[Kolmogorov's exponent is constitutionally forced]
\label{prop:kolmogorov}
The Kolmogorov dissipation exponent $\alpha = 3/4$ is constitutionally forced
at depth $k = 2$.
\end{proposition}

\begin{proof}[Verification against Definition~\ref{def:CF}]
(CF1): The triadic interaction structure of the Navier--Stokes equations gives
a configuration space of size $2^3 = 8$: three independent triadic couplings,
each present or absent.

(CF2): Two constitutional constraints:
\begin{itemize}
  \item \emph{Incompressibility}: $\nabla \cdot \mathbf{u} = 0$. This is
    intrinsic to incompressible fluid dynamics---not a modelling choice.
    It eliminates one configuration class at level $1$.
  \item \emph{Helicity orthogonality}: $\omega^+ \perp \omega^-$. The positive
    and negative helicity components of vorticity are structurally orthogonal,
    forbidding the resonant vorticity coupling necessary for finite-time
    blow-up. This is intrinsic to the vorticity decomposition in three
    dimensions and independent of the incompressibility constraint.
\end{itemize}

(CF3): Invalid configurations: $2$ (one per constitutional constraint). Valid:
$8 - 2 = 6$. Forced constant: $\theta_2 = (2^3 - 2)/2^3 = 6/8 = 3/4$.

(CF4): The exponent $3/4$ is the unique value consistent with both
incompressibility and helicity orthogonality in the triadic interaction space.
Removing either constraint raises the effective exponent; no other value
is geometrically consistent with both.
\end{proof}

\begin{remark}
Proposition~\ref{prop:kolmogorov} provides what 84 years of empirical validation
and dimensional analysis have not: an \textit{a priori} structural reason why
$\alpha = 3/4$ and not any other value. Constitutional Forcing is the reason.
Note also that $k=2$ gives $\theta_2 = 1/2$ in the prime arithmetic context
(the Sophie Germain non-instance), while in the fluid dynamics context $k=2$
acts on a base space of size $2^3 = 8$ rather than $2^2 = 4$, giving $3/4$.
This is not a contradiction: the configuration space structure depends on the
system. The Counting Theorem counts invalid configurations relative to the
appropriate base space for each system.
\end{remark}

%% ─────────────────────────────────────────────────────────────────────────────
\section{Non-Instances and Falsifiability}
\label{sec:noninstances}
%% ─────────────────────────────────────────────────────────────────────────────

A framework that cannot be wrong is not mathematics. Constitutional Forcing
predicts not only which systems exhibit forcing but which do not---and why.

\begin{proposition}[Sophie Germain non-instance]
\label{prop:sophie}
Sophie Germain prime pairs $(p, 2p+1)$ do \emph{not} exhibit Constitutional
Forcing at depth $k = 2$.
\end{proposition}

\begin{proof}
The two apparent modular conditions on Sophie Germain pairs are:
\begin{itemize}
  \item $p \not\equiv 0 \pmod{2}$ (primality of $p$), and
  \item $2p + 1 \not\equiv 0 \pmod{2}$ (primality of $2p+1$).
\end{itemize}
These appear to be two independent constraints (depth $k = 2$, predicting
$\theta_2 = 1/2$). But the second is \emph{not} independent of the first:
if $p$ is an odd prime, then $2p$ is even and $2p+1$ is automatically odd.
The second condition is derivable from the first. Condition (C2) of
Definition~\ref{def:constitutional} fails.

The effective depth is $k = 1$, giving $\theta_1 = 1/2$. This matches the
known level of distribution for Sophie Germain pairs~\cite{WakilShannon}.
The prediction was confirmed computationally before the theoretical verification
was performed.
\end{proof}

\begin{remark}
The Sophie Germain non-instance illustrates a critical distinction: two
constraints that look independent may be coupled, and coupled constraints do
not increase the depth. Constitutional Forcing at depth $k$ requires \emph{genuinely}
independent constraints---independence in the sense of (C2), not merely the
appearance of independence.
\end{remark}

%% ─────────────────────────────────────────────────────────────────────────────
\section{The Constitutional Forcing Conjecture}
\label{sec:conjecture}
%% ─────────────────────────────────────────────────────────────────────────────

\begin{conjecture}[Constitutional Forcing Conjecture]
\label{conj:CF}
Let $S$ be a system exhibiting Constitutional Forcing of depth $k$ in the sense
of Definition~\ref{def:CF}. Then the governing constant of $S$ is uniquely
forced to be
\[
  \thetak = \frac{2^k - k}{2^k},
\]
and this value is achievable: there exist systems exhibiting Constitutional
Forcing at every depth $k \geq 1$.
\end{conjecture}

\begin{remark}
Conjecture~\ref{conj:CF} has two parts. The \emph{uniqueness} part---that the
governing constant is $\thetak$---follows from the Counting Theorem
(Theorem~\ref{thm:count}) whenever (CF1)--(CF3) are verified. The
\emph{achievability} part---that every depth $k$ is realised---is the open
content. The known instances provide $k = 1$ (prime arithmetic and information
theory), $k = 2$ (fluid dynamics), and $k = 3$ (cascade moduli). The first open
case is $k = 4$, predicting a governing constant of $3/4$ in some as-yet
unidentified system with four independent constitutional constraints. The
$k$-hierarchy prediction table (Table~\ref{tab:theta}) is the program's
falsifiable prediction set.
\end{remark}

\begin{remark}
The conjecture connects to the Constitutional Forcing Conjecture as stated in
categorical language in~\cite{WakilShannon}: the category $\mathcal{SW}$ of
Shannon--Wakil systems and the functor $\Phi\colon \mathcal{SW} \to (0,1)$
that assigns the governing constant. Constitutional Forcing is what gives the
functor its value; the Shannon--Wakil Effect is what makes the functor
observable. The two conjectures are the same statement from two perspectives.
\end{remark}

%% ─────────────────────────────────────────────────────────────────────────────
\section{Position in the Seven-Paper Program}
\label{sec:program}
%% ─────────────────────────────────────────────────────────────────────────────

This paper is Paper~00 of a seven-paper program on constitutional sieve theory.
Table~\ref{tab:program} shows how each paper relates to the mechanism defined here.

\begin{table}[h]
\centering
\caption{The seven-paper program. Each paper is an instance or consequence
  of Constitutional Forcing.}
\label{tab:program}
\begin{tabular}{c p{3.8cm} p{5.2cm} c}
\toprule
Paper & Title & Role & CF depth \\
\midrule
00 & Constitutional Forcing (this paper)
   & Mechanism coined & All $k$ \\[2pt]
01 & Khayyam's Triangle~\cite{WakilKhayyam}
   & Oldest known instance ($k{=}3$, mod-3 fractal)
   & $k=3$ \\[2pt]
02 & Spherical Cow Philosophy~\cite{WakilSC}
   & Methodology: finding CF in the wild
   & --- \\[2pt]
03 & Shannon--Wakil Effect~\cite{WakilShannon}
   & Observable phenomenon CF produces
   & $k=1,3$ \\[2pt]
04 & Ramanujan's Dimensional Forcing~\cite{WakilRamanujan}
   & Universal formula $\thetak = (2^k{-}k)/2^k$
   & All $k$ \\[2pt]
05 & Eisenstein Cascade Moduli~\cite{WakilEisenstein}
   & Proof of $\theta_3 = 5/8$ (arithmetic)
   & $k=3$ \\[2pt]
06 & Condition W3~\cite{WakilW3}
   & Siegel zeros excluded; proof closed
   & $k=3$ \\
\bottomrule
\end{tabular}
\end{table}

\noindent The logical dependency structure is:
\[
  \text{Paper~00} \;\longrightarrow\;
  \begin{cases}
    \text{Papers 01, 03, 04} & \text{(instances and cross-domain structure)} \\
    \text{Papers 02} & \text{(methodology)} \\
    \text{Papers 05, 06} & \text{(core arithmetic proof)}
  \end{cases}
\]
A reader wishing to understand the arithmetic proof should read Papers 00, 04,
05, 06 in order. A reader wishing to understand the cross-domain structure
should read Papers 00, 01, 02, 03. A reader who wishes to understand why the
prime $3$ is special should read Papers 00 and 01: Constitutional Forcing
tells you the mechanism; Khayyam's Triangle shows you that it was operating
long before anyone named it.

%% ─────────────────────────────────────────────────────────────────────────────
\section{What Constitutional Forcing Is Not}
\label{sec:not}
%% ─────────────────────────────────────────────────────────────────────────────

\subsection{Not an analogy}

An analogy observes that two things look similar and draws heuristic inspiration.
Constitutional Forcing is not an analogy between number theory and information
theory. It is a common \emph{mechanism} that both instantiate. The difference
is the difference between two fires looking similar and two fires sharing the
same oxygen. Analogies are suggestive; Constitutional Forcing is causal.

\subsection{Not numerology}

The formula $\thetak = (2^k-k)/2^k$ is derived from the Counting Theorem
(Theorem~\ref{thm:count}), which is proved by induction from the definition
of constitutional constraints. The formula is a theorem, not a fit.
Furthermore, the framework is falsifiable: non-instances are predicted before
verification (Proposition~\ref{prop:sophie}), and open predictions at $k = 4$
are testable.

\subsection{Not a claim of isomorphism}

We do not claim that prime arithmetic and information theory are isomorphic, or
that there is a functor between them as mathematical objects. We claim that
both exhibit Constitutional Forcing---that both are instances of the same
mechanism. The same claim could be made about two entirely different physical
systems that both exhibit, say, Hamiltonian dynamics: no isomorphism between
the systems is implied, only a shared structural property.

%% ─────────────────────────────────────────────────────────────────────────────
\section*{Acknowledgements}
%% ─────────────────────────────────────────────────────────────────────────────

We thank Omar Khayyam (1048--1131), whose geometric instinct to ask not
\emph{how many?} but \emph{what shape is hiding here?} produced, in the
mod-3 fractal of his Triangle, the earliest known geometric manifestation
of Constitutional Forcing---nine centuries before anyone named it.

We thank Claude Shannon, whose 1948 theorems contained two instances of
Constitutional Forcing without naming the mechanism, and Andrei Kolmogorov,
whose 1941 scaling exponent contained a third.

Constitutional Forcing was not invented in 2026. It was \emph{recognised}
in 2026. That is a different thing, and the difference matters.

%% ─────────────────────────────────────────────────────────────────────────────
\begin{thebibliography}{99}
%% ─────────────────────────────────────────────────────────────────────────────

\bibitem{BV1965}
E.~Bombieri and A.~I.~Vinogradov,
On the large sieve,
\emph{Mathematika} \textbf{12} (1965), 201--225.

\bibitem{Kolmogorov1941}
A.~N.~Kolmogorov,
The local structure of turbulence in incompressible viscous fluid for very
large Reynolds numbers,
\emph{Dokl.\ Akad.\ Nauk SSSR} \textbf{30} (1941), 301--305.

\bibitem{Pascadi2025}
A.~Pascadi,
Level of distribution $5/8$ for smooth moduli,
arXiv:2505.00653v2, 2025.

\bibitem{Shannon1948}
C.~E.~Shannon,
A mathematical theory of communication,
\emph{Bell System Tech.\ J.} \textbf{27} (1948), 379--423, 623--656.

\bibitem{WakilKhayyam}
K.~Wakil,
Khayyam's triangle: historical restoration and the geometry hidden in the
numbers,
ARC Institute of Knowware, preprint, 2026.

\bibitem{WakilSC}
K.~Wakil,
Spherical cow philosophy: a methodological framework for mathematical
discovery,
ARC Institute of Knowware, preprint, 2026.

\bibitem{WakilShannon}
K.~Wakil,
The Shannon--Wakil effect: constitutional forcing as a universal pattern
in information theory and prime arithmetic,
ARC Institute of Knowware, preprint, 2026.

\bibitem{WakilRamanujan}
K.~Wakil,
Ramanujan's dimensional forcing: a universal formula for sieve levels,
ARC Institute of Knowware, preprint, 2026.

\bibitem{WakilEisenstein}
K.~Wakil,
Ternary forcing on the hexagonal lattice: Eisenstein integers, cascade
moduli, and level of distribution~$5/8$,
ARC Institute of Knowware, preprint, 2026.

\bibitem{WakilW3}
K.~Wakil,
Condition W3: absence of Siegel zeros for characters modulo~$3^K$ via
CM~structure of~$\Z[\omega]$,
ARC Institute of Knowware, preprint, 2026.

\end{thebibliography}

\end{document}
